\chapter{Linear regression, optimization, and regularization}
\label{ch:linear-regression}

\begin{learningobjectives}
After this chapter, you should be able to:
\begin{itemize}
  \item derive least squares from an affine prediction rule and squared loss;
  \item explain projection, identifiability, conditioning, and gradient descent;
  \item distinguish ridge, lasso, and elastic-net regularization; and
  \item evaluate a regression pipeline using residual and out-of-sample evidence.
\end{itemize}
\end{learningobjectives}

\begin{chapterconnection}
Chapter~\ref{ch:supervised-contract} fixed the population, feature cutoff,
target, split, and baseline. We can now study model classes without silently
changing the question. The affine map introduced here will recur in logistic
regression, support-vector machines, and neural networks.
\end{chapterconnection}

\section{The affine predictor}

For $x\in\mathbb{R}^p$, an affine predictor has the form
\[
f_{w,b}(x)=w^Tx+b.
\]
The intercept permits a nonzero prediction at the coordinate origin. Calling
this model ``linear'' is conventional; mathematically it is affine unless the
intercept is absorbed into an augmented feature vector.

With design matrix $X\in\mathbb{R}^{n\times p}$ and target vector
$y\in\mathbb{R}^n$, ordinary least squares minimizes
\[
J(w,b)=\frac{1}{n}\sum_{i=1}^n(y_i-w^Tx_i-b)^2.
\]
Squared loss penalizes large residuals strongly and targets the conditional mean
under the population risk. That property is useful only when the mean and the
cost implied by squared error match the task.

\section{Projection and identifiability}

After including an intercept column, write the design matrix as $A$. The fitted
vector $A\widehat\theta$ is the orthogonal projection of $y$ onto the column
space of $A$. At an optimum the residual $r=y-A\widehat\theta$ satisfies
$A^Tr=0$. This geometric statement remains meaningful when coefficients are
not unique.

If $A$ has full column rank, the normal equations give
\[
\widehat\theta=(A^TA)^{-1}A^Ty.
\]
Software should not explicitly form the inverse. QR or singular-value methods
are generally more stable, while iterative methods are useful at larger scale.
When columns are linearly dependent, multiple coefficient vectors produce the
same fitted values. Prediction may remain defined while individual coefficient
interpretations do not.

\conceptcheck{Two columns record temperature in Celsius and Fahrenheit. What
happens to rank, fitted values, and the uniqueness of their separate
coefficients when an intercept is present?}

\section{Gradient descent exposes the optimization problem}

For centered notation without an explicit intercept,
\[
\nabla J(w)=\frac{2}{n}X^T(Xw-y).
\]
Gradient descent updates $w_{t+1}=w_t-\eta_t\nabla J(w_t)$. The learning rate
$\eta_t$ converts a direction into a step. A rate that is too large can diverge;
a very small rate wastes computation. Feature scales change the curvature and
therefore the usable step size.

Standardization can improve conditioning, but it is a fitted transformation.
Its means and scales belong to the training fold and serialized pipeline.
Convergence of the numerical optimizer says that an objective was minimized; it
does not say that the model generalizes or that the objective represents the
scientific goal.

\begin{professionalbox}
Record the objective definition, initialization, seed where relevant, stopping
criterion, iteration count, numerical precision, and convergence status. A
training function should not silently return the last iterate after divergence
or a nonfinite loss.
\end{professionalbox}

\section{Basis expansion changes the feature map}

Polynomial regression remains linear in its fitted coefficients. A map
$\phi(x)=(1,x,x^2,\ldots,x^d)$ creates a richer function
$f(x)=w^T\phi(x)$. Interactions can represent joint effects. The flexibility
comes from the basis, not from abandoning linear estimation.

High-degree bases can be poorly conditioned and unstable outside the observed
range. Every basis term must be generated identically during training and
inference. Feature names, order, units, and transformation parameters therefore
belong to the artifact contract.

\section{Regularization chooses among fitting rules}

Ridge regression minimizes
\[
\widehat R(w)+\lambda\lVert w\rVert_2^2,
\]
shrinking correlated or weakly identified coefficients and usually improving
conditioning. Lasso uses $\lambda\lVert w\rVert_1$ and can set coefficients
exactly to zero. Elastic net combines both penalties.

Regularization trades empirical fit for a restriction on the solution. The
penalty depends on feature scale, so scaling and whether the intercept is
penalized are part of the definition. A lasso-selected variable set is not an
automatic list of scientific causes; selection can be unstable among correlated
features.

\begin{misconceptionbox}
\textbf{Misconception:} regularization is needed only when a model visibly
overfits. \textbf{Correction:} it can also stabilize an ill-conditioned
problem, encode a preference among many solutions, and improve performance
under finite samples. Its strength must still be chosen inside evaluation.
\end{misconceptionbox}

\section{Residuals are structured evidence}

A residual $e_i=y_i-\widehat y_i$ retains sign, scale, and context. Plot it
against predictions, major features, time, and groups. Curvature can reveal
missing structure; a fan shape can reveal changing conditional variance;
clusters can reveal omitted batches; isolated points can dominate squared loss.

Aggregate error should include units and a distribution, not only a mean.
Compare mean absolute error, root mean squared error, and relevant quantiles.
Inspect error by scientifically meaningful subgroup. Prediction intervals,
confidence intervals for parameters, and variability across resampled training
sets answer different questions and should not be conflated.

\begin{workedexample}{regularized capacity prediction}
Suppose early-cycle summaries include mean voltage, resistance change, and two
strongly correlated temperature summaries. A least-squares fit produces large
opposite temperature coefficients while predictions change little. This is a
signature of weak coefficient identification.

Fit a pipeline that learns scaling and ridge regression using only each
training fold. Choose $\lambda$ on validation batches, then compare once on the
held-out later batches. Report the baseline and ridge errors in ampere-hours,
residuals by batch, and the range of each input. The ridge coefficients are more
stable, but the report avoids interpreting their magnitudes as causal effects.
\end{workedexample}

\section{Exercises}
\begin{enumerate}
  \item Derive the least-squares gradient with an intercept and verify its
  dimensions.
  \item Construct a rank-deficient design matrix with unique fitted values but
  nonunique coefficients. Explain how ridge changes the solution.
  \item Show why fitting a scaler before cross-validation leaks information.
  \item Design tests for a regression artifact when columns are reordered,
  missing, duplicated, or expressed in incompatible units.
\end{enumerate}

\begin{chapterreview}
Linear regression combines an affine representation, a loss, an optimization
method, and a data contract. Projection explains fitted values; rank and
conditioning explain coefficient instability; gradient descent separates
optimization from estimation; and regularization selects a more constrained
solution. Out-of-sample residual evidence, not training fit alone, supports the
prediction claim.
\end{chapterreview}

\begin{notebooklink}
Lecture 10 notebook 00 develops affine models and least squares. Later notebooks
compare gradient-based optimization, ridge, lasso, diagnostics, and a packaged
regression pipeline.
\end{notebooklink}
